%======================================================================
\section{Two sharp rates}\label{sec:rates}
%======================================================================

The state-weighted method of \cref{sec:state,sec:count,sec:transfer} is not specific to lamp groups. This section applies it to two games
for which the earlier lower bounds lose a power of $\eps$ on the way from the strategy to a count. For the XOR games CHSH$(n)$ the qubit count of \cref{thm:count}, applied once in the reduced state
of the players, gives the ebit cost $\Theta(\min(n,1/\eps))$ (Theorem~E, \cref{sec:rates-chsh}). For the linear system game of Slofstra
and Vidick, a transport argument for spectral measures in one state gives the profile $\Theta(1/\eps)$ of its group and, through
\cref{thm:transfer}, the local dimension $\Theta(\eps^{-1/2})$ (Theorem~F, \cref{sec:rates-sv}). In both cases the new lower bound meets an
upper bound of the same order: for CHSH$(n)$ a known one, and for the game of Slofstra and Vidick theirs, through a lift that keeps the
distinguished involution balanced (\cref{lem:sv-lift}).

\subsection{The CHSH\texorpdfstring{$(n)$}{(n)} games}\label{sec:rates-chsh}

The XOR games CHSH$(n)$ of Slofstra~\cite{Slo11} force $n$ pairwise anticommuting observables on one side. A strategy within a factor
$1-\eps$ of the optimal bias pays a CHSH deficit for each of the $\binom n2$ pairs, and these deficits add up to $O(n^2\eps)$. From triples
of anticommuting observables one builds $\Theta(n)$ approximate qubits~\cite{OV18}, and \cref{thm:count} charges each of them only the sum
of the squares of its pair errors, which is $O(n\eps)$. So a constant fraction of the $n/3$ qubits survives as long as $n\eps$ is small.

\begin{definition}[CHSH$(n)$; \cite{Slo11}, {\cite[Section~2.3]{OV18}}]\label{def:chsh}
In CHSH$(n)$, $n\ge2$, the referee picks $i<j$ uniformly, sends Alice $i$ or $j$ and Bob $(i,j)$ or $(j,i)$, independently and uniformly,
and accepts if and only if $a\oplus b=1$ when Alice received $j$ and Bob $(j,i)$, and $a\oplus b=0$ otherwise. A strategy is a unit vector
$\psi\in H_A\otimes H_B$, with $H_A,H_B$ finite-dimensional, together with binary observables (Hermitian contractions) $A_k$ for Alice,
$k\in[n]$, and $B_{ij}$ for Bob, $i\ne j$. Its \emph{bias}, the probability of winning minus that of losing, is
\[
\operatorname{bias}=\frac1{2n(n-1)}\sum_{i<j}C_{ij},\qquad C_{ij}=\ip{\psi}{(A_i\otimes(B_{ij}+B_{ji}))\psi}+\ip{\psi}{(A_j\otimes(B_{ij}-B_{ji}))\psi}.
\]
The optimal bias is $\operatorname{bias}^*=\sqrt2/2$, and a strategy is \emph{$\eps$-optimal} if $\operatorname{bias}\ge(1-\eps)\operatorname{bias}^*$.
\end{definition}

Ostrev and Vidick measure the bias by its absolute value; replacing every $A_k$ by $-A_k$ changes its sign and nothing else, so we may take
it nonnegative. The next lemma is the robust form of Tsirelson's bound for one pair, with every error measured in Alice's reduced state.

\begin{lemma}[One pair]\label{lem:one-pair}
Let $A,A'$ be involutions on $H_A$, $B,B'$ involutions on $H_B$, $\psi\in H_A\otimes H_B$ a unit vector and $\xi_A=\rho_A^{1/2}$. Put
$T_0=(B+B')/\sqrt2$, $T_1=(B-B')/\sqrt2$,
$C=\ip{\psi}{(A\otimes(B+B'))\psi}+\ip{\psi}{(A'\otimes(B-B'))\psi}$ and $\delta=2\sqrt2-C$. Then
\[
\delta=\frac1{\sqrt2}\Bigl(\lVert(A\otimes I-I\otimes T_0)\psi\rVert^2+\lVert(A'\otimes I-I\otimes T_1)\psi\rVert^2\Bigr).
\]
Consequently $\delta\ge0$, and
\[
\hs{\{A,A'\}\xi_A}^2\le33\,\delta,\qquad\hs{[A,\xi_A]}^2\le2\sqrt2\,\delta,\qquad\hs{[A',\xi_A]}^2\le2\sqrt2\,\delta.
\]
\end{lemma}

\begin{proof}
Write $\langle\cdot\rangle$ for the expectation in $\psi$. Since $A^2=I$, $\lVert(A\otimes I-I\otimes T_0)\psi\rVert^2=1+\langle
T_0^2\rangle-\sqrt2\langle A\otimes(B+B')\rangle$, and likewise for the second norm. As $T_0^2+T_1^2=B^2+B'^2=2I$, the sum of the two norms
is $4-\sqrt2\,C$, which is the identity. So both norms are at most $\tau=(\sqrt2\delta)^{1/2}$. Since $\{T_0,T_1\}=B^2-B'^2=0$ and
$\op{T_0},\op{T_1}\le\sqrt2$, and since $A\otimes I$ commutes with $I\otimes T_1$,
\[
(AA'\otimes I)\psi=(A\otimes I)(I\otimes T_1)\psi+(A\otimes I)\bigl((A'\otimes I)-(I\otimes T_1)\bigr)\psi=(I\otimes T_1T_0)\psi+\epsilon_1,
\]
with $\lVert\epsilon_1\rVert\le(1+\sqrt2)\tau$. In the same way $(A'A\otimes I)\psi=(I\otimes T_0T_1)\psi+\epsilon_2$ with
$\lVert\epsilon_2\rVert\le(1+\sqrt2)\tau$. Adding, and using \cref{lem:standard-form}(a) for Alice's side,
$\hs{\{A,A'\}\xi_A}\le2(1+\sqrt2)\tau$; and $4(1+\sqrt2)^2\sqrt2<33$. Finally \cref{lem:standard-form}(b), with the roles of the two
factors exchanged, gives Hermitian operators $\tilde T_0,\tilde T_1$ on $H_A$ with $\hs{A\xi_A-\xi_A\tilde T_0}$,
$\hs{A'\xi_A-\xi_A\tilde T_1}\le\tau$, and \cref{lem:consistency} gives the last two bounds.
\end{proof}

With all observables involutions, the lemma gives $C_{ij}\le2\sqrt2$, which is Tsirelson's bound~\cite{Cir80}, and hence
$\operatorname{bias}\le\operatorname{bias}^*$; the reduction to involutions in the next proof shows that this holds for all strategies.

\begin{theorem}[Ebits for CHSH$(n)$]\label{thm:chsh}
Let $n\ge3$ and $\eps>0$, and put $M=\lfloor\min(n,10^{-7}/\eps)/3\rfloor$. Every $\eps$-optimal strategy for CHSH$(n)$ satisfies
\[
S(\rho_A)\ \ge\ M/2 .
\]
In particular $S(\rho_A)\ge10^{-7}\min(n,1/\eps)/6-1/2$.
\end{theorem}

The proof selects $n'=3M$ questions whose pairwise CHSH deficits are small on average, splits them into $M$ triples, turns each triple
into one approximate Pauli pair as in~\cite{OV18}, and applies \cref{thm:count} to the $M$ pairs in Alice's reduced state. The pair errors
enter only through their squares, which add up to $O(n'^2\eps)$, that is, $O(n'\eps)$ per qubit.

\begin{proof}
If $M=0$ there is nothing to prove, so let $M\ge1$ and $n'=3M\le\min(n,10^{-7}/\eps)$.

\emph{Step 0: reductions.} With all other observables fixed, the bias is an affine function $\Tr(A_kN_k)+\mathrm{const}$ of Alice's
observable $A_k$, with $N_k$ Hermitian, so replacing $A_k$ by $\sgn(N_k)$ (with $\sgn(0)=1$) does not lower it; the same holds for Bob's
observables. Doing this for each observable in turn, we may assume that all observables are involutions; the state is unchanged. Let
$\delta_{ij}=2\sqrt2-C_{ij}\ge0$ (\cref{lem:one-pair}). From $\operatorname{bias}\ge(1-\eps)\sqrt2/2$,
\[
\sum_{i<j}\delta_{ij}\ \le\ 2\sqrt2\,\eps\binom n2 .
\]
Averaging over uniformly random $n'$-subsets of $[n]$, fix a subset $\mathcal N$ whose pair average
$\bar\delta=\binom{n'}2^{-1}\sum_{x<y\in\mathcal N}\delta_{xy}$ is at most $2\sqrt2\,\eps$. A fixed pair of $\mathcal N$ lies in a common block of a
uniformly random partition of $\mathcal N$ into $M$ triples with probability $2/(n'-1)$; fix a partition $\mathcal T_1,\dots,\mathcal T_M$
for which the sum of $\delta_{xy}$ over the $3M$ pairs inside the triples is at most $\frac2{n'-1}\binom{n'}2\bar\delta=n'\bar\delta$.

\emph{Step 1: pair errors.} Write $\xi=\rho_A^{1/2}$, $\eta_{xy}=\hs{\{A_x,A_y\}\xi}$ and $c_x=\hs{[A_x,\xi]}$. Applying \cref{lem:one-pair}
to the pair $\{x,y\}$, with $(A,A',B,B')=(A_x,A_y,B_{xy},B_{yx})$ for $x<y$,
\[
\eta_{xy}^2\le33\,\delta_{xy},\qquad c_x^2\le2\sqrt2\,\delta_{xy}\quad\text{for every }y\ne x.
\]
Choosing for each $x\in\mathcal N$ the partner $y\in\mathcal N$ with the least $\delta_{xy}$,
\[
\sum_{x\in\mathcal N}c_x^2\ \le\ 2\sqrt2\sum_{x\in\mathcal N}\frac1{n'-1}\sum_{y\in\mathcal N\setminus x}\delta_{xy}\ =\ 2\sqrt2\,n'\bar\delta .
\]

\emph{Step 2: one qubit from each triple.} The construction of pairs from triples is that of~\cite{OV18}. For $\mathcal T_k=\{a,b,c\}$,
labelled in any order, put $U_k=iA_aA_b$ and $W_k=iA_bA_c$, which are unitaries, and let $X_k=\mathrm r(U_k)$ and $Z_k=\mathrm r(W_k)$ be
their roundings (\cref{lem:rounding}). Since $U_k^2-I=-A_aA_b\{A_a,A_b\}$, we have $\hs{(U_k^2-I)\xi}=\eta_{ab}$, so
\[
\hs{(X_k-U_k)\xi}\le\eta_{ab},\qquad\hs{[X_k,\xi]}\le c_a+c_b+2\eta_{ab},
\]
and likewise for $Z_k$ with $(b,c)$ in place of $(a,b)$. Put $\lambda_k=\eta_{ab}+\eta_{bc}+\eta_{ac}+c_a+c_b+c_c$ and, for $k\ne l$,
$\mu_{kl}=\sum_{x\in\mathcal T_k,\,y\in\mathcal T_l}\eta_{xy}$. For distinct indices $b,c,b',c'$, and for one triple $\{a,b,c\}$, we have
the identities, which use only $A_x^2=I$,
\begin{align*}
A_bA_cA_{b'}A_{c'}-A_{b'}A_{c'}A_bA_c&=A_b\{A_c,A_{b'}\}A_{c'}-\{A_b,A_{b'}\}A_cA_{c'}\\
&\quad+A_{b'}A_b\{A_c,A_{c'}\}-A_{b'}\{A_b,A_{c'}\}A_c,\\
\{U_k,W_k\}&=-A_b\{A_a,A_c\}A_b+\{A_a,A_b\}A_cA_b-A_aA_b\{A_b,A_c\}.
\end{align*}
Apply them to $\xi$ and move the trailing observables to the right of $\xi$, at cost $2c_x$ for each, the anticommutators having norm at
most $2$. For $k\ne l$ this gives
\[
\hs{[W_k,W_l]\xi},\ \hs{[U_k,W_l]\xi}\ \le\ \mu_{kl}+4\lambda_k+4\lambda_l,\qquad\hs{\{U_k,W_k\}\xi}\le4\lambda_k .
\]
\Cref{lem:calculus}(e), with the primed operators the unitaries $U$ and $W$, whose state commutators are at most $\lambda_k$, replaces
$U,W$ by $X,Z$: for $k\ne l$,
\begin{gather*}
\hs{[Z_k,Z_l]\xi},\ \hs{[X_k,Z_l]\xi}\le\mu_{kl}+8\lambda_k+8\lambda_l,\\
\hs{\{X_k,Z_k\}\xi}\le12\lambda_k,\qquad\hs{[X_k,\xi]},\ \hs{[Z_k,\xi]}\le2\lambda_k .
\end{gather*}

\emph{Step 3: counting.} Apply \cref{thm:count} to the $M$ pairs $(X_k,Z_k)$ in the state $\rho_A$. With
$c^X_k=\hs{[X_k,\xi]}$ and $c^Z_k=\hs{[Z_k,\xi]}$, its parts (i)--(iii) give numbers $a_k,b_k\ge0$ with
$a_k^2\le2\sum_{l>k}(\hs{[Z_k,Z_l]\xi}+c^Z_l)^2+2\sum_{l>k}(c^Z_l)^2$,
$b_k^2\le2(\hs{\{X_k,Z_k\}\xi}+c^Z_k)^2+2\sum_{l\ne k}(\hs{[X_k,Z_l]\xi}+c^Z_l)^2+2\sum_{l}(c^Z_l)^2$ and
\[
S(\rho_A)\ \ge\ \sum_{k=1}^M\Bigl[1-\frac{(b_k+c^X_k)^2}{2\ln2}-h\Bigl(\min\Bigl(\frac{(a_k+c^Z_k)^2}4,\frac12\Bigr)\Bigr)\Bigr].
\]
Put $\Sigma_\lambda=\sum_k\lambda_k^2$. By Step~2 and $(x+y+z)^2\le3(x^2+y^2+z^2)$,
\[
a_k^2\le6\sum_{l\ne k}\mu_{kl}^2+384M\lambda_k^2+608\,\Sigma_\lambda,\qquad b_k^2\le392\lambda_k^2+6\sum_{l\ne k}\mu_{kl}^2+384M\lambda_k^2+608\,\Sigma_\lambda .
\]
By $(x+y)^2\le2x^2+2y^2$, the monotonicity of $h$ on $[0,\frac12]$, and Jensen's inequality for the concave nondecreasing function
$t\mapsto h(\min(t,\frac12))$,
\[
S(\rho_A)\ \ge\ M-\frac{B}{\ln2}-M\,h\Bigl(\min\Bigl(\frac A{2M},\frac12\Bigr)\Bigr),\qquad A=\sum_k(a_k^2+4\lambda_k^2),\quad B=\sum_k(b_k^2+4\lambda_k^2).
\]
Now $\mu_{kl}^2\le9\sum_{x\in\mathcal T_k,y\in\mathcal T_l}\eta_{xy}^2$, so $\sum_{k\ne l}\mu_{kl}^2\le18\sum_{x<y\in\mathcal
N}\eta_{xy}^2\le594\binom{n'}2\bar\delta\le297\,n'^2\bar\delta$, and $\Sigma_\lambda\le6\bigl(\sum_{\rm internal}\eta_{xy}^2+\sum_{x\in\mathcal
N}c_x^2\bigr)\le6(33+2\sqrt2)\,n'\bar\delta\le215\,n'\bar\delta$, where the internal pairs are those inside the triples. Hence, with
$M\ge1$ and $M=n'/3$,
\[
A\le6\cdot297\,n'^2\bar\delta+996\,M\Sigma_\lambda\le73162\,n'^2\bar\delta,\qquad B\le6\cdot297\,n'^2\bar\delta+1388\,M\Sigma_\lambda\le101256\,n'^2\bar\delta .
\]
With $\bar\delta\le2\sqrt2\,\eps$ this gives
\[
S(\rho_A)\ \ge\ M\Bigl[1-1.24\cdot10^6\,n'\eps-h\bigl(\min(3.11\cdot10^5\,n'\eps,\tfrac12)\bigr)\Bigr].
\]
Since $n'\eps\le10^{-7}$, the bracket is at least $1-0.124-h(0.0311)>0.67$, so $S(\rho_A)\ge M/2$. Finally
$M\ge\min(n,10^{-7}/\eps)/3-1\ge10^{-7}\min(n,1/\eps)/3-1$.
\end{proof}

\begin{corollary}[The ebit cost of CHSH$(n)$]\label{cor:chsh-ebits}
Let $S_{\min}(n,\eps)$ be the least entanglement entropy $S(\rho_A)$ of an $\eps$-optimal strategy for CHSH$(n)$. For $n\ge3$ and
$0<\eps\le1$,
\[
\frac{10^{-7}}6\min(n,1/\eps)-\frac12\ \le\ S_{\min}(n,\eps)\ \le\ \min\bigl(\lfloor n/2\rfloor,\lceil1/\eps\rceil\bigr).
\]
So $S_{\min}(n,\eps)=\Theta(\min(n,1/\eps))$ as $\min(n,1/\eps)\to\infty$.
\end{corollary}

\begin{proof}
The lower bound is \cref{thm:chsh}. For the upper bound $\lceil1/\eps\rceil$, every XOR game has an $\eps$-optimal strategy on a maximally
entangled state of local dimension $2^{\lceil1/\eps\rceil}$~\cite[Theorem~2]{OV18}. For the upper bound $\lfloor n/2\rfloor$, let
$m=\lfloor n/2\rfloor$; the Jordan--Wigner construction gives $2m+1\ge n$ pairwise anticommuting involutions $A_1,\dots,A_n$ on
$(\C^2)^{\otimes m}$. On the maximally entangled state $\Phi$ of local dimension $2^m$, for which $\ip{\Phi}{(Y\otimes Z^T)\Phi}=2^{-m}\Tr(YZ)$,
let Bob use $B_{ij}=((A_i+A_j)/\sqrt2)^T$ and $B_{ji}=((A_i-A_j)/\sqrt2)^T$ for $i<j$, which are involutions. Then
$B_{ij}+B_{ji}=\sqrt2A_i^T$ and $B_{ij}-B_{ji}=\sqrt2A_j^T$, so $C_{ij}=2\sqrt2$ for every pair and
$\operatorname{bias}=\operatorname{bias}^*$, with $m$ ebits~\cite{Slo11,OV18}.
\end{proof}

The entropy bound does not use purity in an essential way: shared randomness cannot lower the entanglement of formation.

\begin{corollary}[Entanglement of formation]\label{cor:chsh-ef}
Let $\sigma$ be a state on $H_A\otimes H_B$ and let binary observables for CHSH$(n)$, $n\ge3$, achieve bias at least
$(1-\eps)\operatorname{bias}^*$ on $\sigma$. Then
\[
E_F(\sigma)\ \ge\ \frac14\Bigl\lfloor\frac13\min\Bigl(n,\frac{10^{-7}}{2\eps}\Bigr)\Bigr\rfloor .
\]
\end{corollary}

\begin{proof}
Let $\sigma=\sum_jp_j\lvert\psi_j\rangle\langle\psi_j\rvert$ be any decomposition into pure states. The bias is linear in the state, so
$\operatorname{bias}(\sigma)=\sum_jp_j\operatorname{bias}(\psi_j)$ with the same observables, and
$\operatorname{bias}(\psi_j)\le\operatorname{bias}^*$. Put $\eps_j=1-\operatorname{bias}(\psi_j)/\operatorname{bias}^*\ge0$; then
$\sum_jp_j\eps_j\le\eps$, so the components with $\eps_j\le2\eps$ have total weight at least $\frac12$. Each of them is a $2\eps$-optimal
pure strategy, so its reduced state on $H_A$ has entropy at least half the floor in the statement, by \cref{thm:chsh}. Averaging over $j$,
with entropy $\ge0$ for the other components, and taking the infimum over decompositions gives the claim.
\end{proof}

The strategies of \cref{cor:chsh-ebits} use pure states, for which $E_F$ is the entanglement entropy, so the entanglement of formation
needed is also $\Theta(\min(n,1/\eps))$. \emph{Theorem~E is \cref{thm:chsh} with \cref{cor:chsh-ebits,cor:chsh-ef}.}

\begin{remark}\label{rates:rem-chsh}
The rate is linear in $1/\eps$, which~\cite[Theorem~2]{OV18} shows is optimal; the lower bound of~\cite[Theorem~1]{OV18} is
$\Omega(\eps^{-1/5})$, with the same question sets of sizes $n$ and $n(n-1)$. The argument of~\cite{OV18} also builds about $n/3$ pairs from
triples, but converts approximate anticommutation into a qubit count with polynomial losses at each step. Here every error stays in the
state-weighted $\ell^2$ norm, and \cref{thm:count} is applied once. The threshold $n\eps\approx1$ at which the count stops is the scale at
which the maximally entangled strategy of~\cite[Theorem~2]{OV18}, with $\lceil1/\eps\rceil$ ebits, becomes cheaper than the exact strategy
with $\lfloor n/2\rfloor$.
\end{remark}

\subsection{The Slofstra--Vidick group and game}\label{sec:rates-sv}

Slofstra and Vidick built a linear system game $\Gsv$ with $C\eps^{-1/6}\le\Ent(\Gsv,\eps)\le C'\eps^{-1/2}$~\cite[Theorem~1.1]{SV18} from
the group~\cite[Proposition~3.2]{SV18}
\begin{equation}\label{eq:K}
K=\bigl\langle a,b,c,x,y\ \bigm|\ xyx^{-1}y^{-2},\ xcx^{-1}c^{-1},\ yay^{-1}b^{-1},\ yby^{-1}a^{-1},\ c^{-1}ab,\ a^2,\ b^2,\ c^2\bigr\rangle,
\end{equation}
in which $c$ is a central involution, and whose profile satisfies $C\eps^{-2/3}\le\hlp_K(c;\delta,\eps)\le C'/\eps$. We show that both
lower exponents are the upper ones: $\hlp_K(c;\delta,\eps)=\Theta(1/\eps)$ (\cref{cor:K-profile}) and $\Ent(\Gsv,\eps)=\Theta(\eps^{-1/2})$
(\cref{thm:sv-game}).

The picture is as follows. Let $\mu$ be the spectral measure of $\phi(y)$ in the state. The relation $xyx^{-1}=y^2$ makes $\mu$ nearly
invariant under the doubling map $\vartheta\mapsto2\vartheta$ of the circle, and, where $c$ acts as $-1$, the relations $yay^{-1}=b$ and
$c=ab$ make it nearly invariant under the half-turn $\vartheta\mapsto\vartheta+\pi$. Averaging over the two inverse branches of the
doubling map contracts the Wasserstein distance by a factor $2$, so a measure that is nearly invariant under both maps is close to Haar
measure (\cref{prop:K-W1}), while a measure with $d$ atoms is at distance at least $\pi/(2d)$ from it. All quantities are measured in
one state; the
Wasserstein-$1$ distance, rather than a quadratic one, is what makes a change of state cost only a linear error.

\subsubsection*{Measures on the circle}

Let $\T=\R/2\pi\Z$ with the geodesic distance $d(\vartheta,\vartheta')\in[0,\pi]$, and let $\mathsf h$ be the normalized Haar measure. For
probability measures $\mu,\nu$ on $\T$ put
\[
W_1(\mu,\nu)=\sup\Bigl\{\int f\,d\mu-\int f\,d\nu\ :\ f:\T\to\R\text{ is $1$-Lipschitz}\Bigr\}.
\]
By Kantorovich--Rubinstein duality~\cite[Chapter~5]{Vil09} this is the least mean distance of a coupling of $\mu$ and $\nu$; we use only
the triangle inequality and the elementary bound $W_1(\mu,\nu)\le\int d\,d\gamma$ for every coupling $\gamma$ of $\mu$ and $\nu$, which
holds because $\int f\,d\mu-\int f\,d\nu=\int(f(\vartheta)-f(\vartheta'))\,d\gamma$. Let $\mathsf D(\vartheta)=2\vartheta$ and
$\mathsf S(\vartheta)=\vartheta+\pi$, with push-forwards $\mathsf D_*$, $\mathsf S_*$. The two square roots of $e^{i\vartheta}$ have angles
$\vartheta/2$ and $\vartheta/2+\pi$, a set that does not depend on the representative of $\vartheta$. Let $\mathsf P$ replace each point by
the uniform measure on its two square roots: $\int f\,d(\mathsf P\mu)=\int\mathsf P^*f\,d\mu$ with
$\mathsf P^*f(\vartheta)=\frac12\bigl(f(\vartheta/2)+f(\vartheta/2+\pi)\bigr)$.

\begin{lemma}[Doubling and half-turn]\label{lem:circle}
For probability measures $\mu,\nu$ on $\T$:
\begin{enumerate}[label=(\alph*),itemsep=1pt]
\item $\mathsf P\mathsf h=\mathsf h$ and $\mathsf P(\mathsf D_*\mu)=\frac12(\mu+\mathsf S_*\mu)$;
\item $W_1(\mathsf P\mu,\mathsf P\nu)\le\frac12W_1(\mu,\nu)$;
\item $W_1(\mu,\mathsf h)\le W_1(\mathsf D_*\mu,\mu)+W_1(\mathsf S_*\mu,\mu)$;
\item if $\mu$ is supported on at most $d$ points, then $W_1(\mu,\mathsf h)\ge\pi/(2d)$.
\end{enumerate}
\end{lemma}

\begin{proof}
(a) The maps $\vartheta\mapsto\vartheta/2$ and $\vartheta\mapsto\vartheta/2+\pi$, on a fundamental interval of length $2\pi$, carry $\mathsf h$ to
the uniform measures on two complementary half-circles, so $\int\mathsf P^*f\,d\mathsf h=\int f\,d\mathsf h$. The square roots of
$e^{2i\vartheta}$ are $\vartheta$ and $\vartheta+\pi$.

(b) Let $f$ be $1$-Lipschitz and $d(\vartheta,\vartheta')=t$. Choose real representatives with $\vartheta'=\vartheta+s$ and $\lvert
s\rvert=t\le\pi$. The roots $\vartheta/2$ and $(\vartheta+s)/2$, and the roots $\vartheta/2+\pi$ and $(\vartheta+s)/2+\pi$, are at distance at
most $\lvert s\rvert/2$, so $\lvert\mathsf P^*f(\vartheta)-\mathsf P^*f(\vartheta')\rvert\le t/2$. Thus $2\mathsf P^*f$ is $1$-Lipschitz and
$\int f\,d(\mathsf P\mu-\mathsf P\nu)=\int\mathsf P^*f\,d(\mu-\nu)\le\frac12W_1(\mu,\nu)$.

(c) By the triangle inequality, (a) and (b),
\begin{align*}
W_1(\mu,\mathsf h)&\le W_1(\mu,\mathsf P\mathsf D_*\mu)+W_1(\mathsf P\mathsf D_*\mu,\mathsf P\mu)+W_1(\mathsf P\mu,\mathsf P\mathsf h)\\
&\le\tfrac12W_1(\mu,\mathsf S_*\mu)+\tfrac12W_1(\mathsf D_*\mu,\mu)+\tfrac12W_1(\mu,\mathsf h),
\end{align*}
where the first term was computed directly: $\int f\,d(\mu-\frac12(\mu+\mathsf S_*\mu))=\frac12\int f\,d(\mu-\mathsf S_*\mu)$. Since
$W_1(\mu,\mathsf h)\le\pi$ is finite, (c) follows.

(d) Let $k\le d$ be the number of points of $\supp\mu$. They cut $\T$ into $k$ arcs of lengths $l_1,\dots,l_k$ with $\sum_jl_j=2\pi$. The function
$f(\vartheta)=d(\vartheta,\supp\mu)$ is $1$-Lipschitz, $\int f\,d\mu=0$, and, by Cauchy--Schwarz,
\[
\int f\,d\mathsf h=\frac1{2\pi}\sum_j\int_0^{l_j}\min(t,l_j-t)\,dt=\sum_j\frac{l_j^2}{8\pi}\ \ge\ \frac{(2\pi)^2}{8\pi k}=\frac\pi{2k}. \qedhere
\]
\end{proof}

The operator $\mathsf P$ is the transfer operator of the doubling map, and (b) is an averaged form of the square-root lifting that Slofstra
and Vidick use in operator norm~\cite[Claims~3.6 and~3.7]{SV18}. Part~(d) needs no lower bound on the masses of the atoms.

For a unitary $Y$ and a density matrix $\rho$ on $\C^d$ let $\mu^\rho_Y$ be the \emph{spectral measure of $Y$ in $\rho$}, the probability
measure on $\T$ with $\int f\,d\mu^\rho_Y=\Tr(\rho f(Y))$, where $f(Y)$ applies $f$ to the angles of the eigenvalues. It is supported on
the angles of the at most $d$ eigenvalues of $Y$, and
\begin{equation}\label{eq:rates-spectral}
\mu^\rho_{Y^2}=\mathsf D_*\mu^\rho_Y,\qquad\mu^\rho_{-Y}=\mathsf S_*\mu^\rho_Y,\qquad\mu^\rho_{UYU^*}=\mu^{U^*\rho U}_Y\quad(U\text{ unitary}).
\end{equation}

\begin{lemma}[Coupling in one state]\label{lem:coupling}
Let $P,Q,Y,U$ be unitaries on $\C^d$ and $\xi=\rho^{1/2}$.
\begin{enumerate}[label=(\alph*),itemsep=1pt]
\item There is a coupling $\nu$ of $\mu^\rho_P$ and $\mu^\rho_Q$, viewed on the unit circle, with $\int\lvert z-z'\rvert^2\,d\nu(z,z')=\hs{P\xi-\xi
Q}^2$. Hence $W_1(\mu^\rho_P,\mu^\rho_Q)\le\frac\pi2\hs{P\xi-\xi Q}$.
\item $W_1(\mu^{U^*\rho U}_Y,\mu^\rho_Y)\le\pi\hs{[U,\xi]}$.
\end{enumerate}
\end{lemma}

\begin{proof}
(a) Let $P=\sum_iz_iE_i$ and $Q=\sum_jz'_jF_j$ be spectral decompositions with distinct eigenvalues, and
$\nu=\sum_{i,j}\hs{E_i\xi F_j}^2\,\delta_{(z_i,z'_j)}$. The operators $E_i\xi F_j$ are mutually orthogonal in the Hilbert--Schmidt inner
product. Since $\sum_jF_j=I$, $\sum_j\hs{E_i\xi F_j}^2=\hs{E_i\xi}^2=\Tr(E_i\rho)$, and likewise $\sum_i\hs{E_i\xi F_j}^2=\Tr(\rho F_j)$, so
$\nu$ couples $\mu^\rho_P$ and $\mu^\rho_Q$; and $\int\lvert z-z'\rvert^2d\nu=\hs{\sum_{i,j}(z_i-z'_j)E_i\xi F_j}^2=\hs{P\xi-\xi Q}^2$. The arc
length between two points of the unit circle is at most $\frac\pi2$ times their chordal distance, and $\int\lvert z-z'\rvert\,d\nu\le(\int\lvert
z-z'\rvert^2d\nu)^{1/2}$.

(b) Let $f$ be $1$-Lipschitz and $f_0=f-c_f$, with $c_f$ the midpoint of the range of $f$, so that $\lvert f_0\rvert\le\pi/2$, the diameter
of $\T$ being $\pi$. Both measures have mass one, so $\int f\,d(\mu^{U^*\rho U}_Y-\mu^\rho_Y)=\Tr\bigl((U^*\rho U-\rho)f_0(Y)\bigr)\le\frac\pi2\tn{U^*\rho
U-\rho}$, which is at most $\pi\hs{[U,\xi]}$ by \cref{lem:rotate}.
\end{proof}

The coupling in (a) is Connes' joint distribution trick~\cite[Lemma~1.2.6]{Con76}, in the form used in~\cite[Lemma~5.5]{SV18}; for
$\rho=I/d$ it gives the Hoffman--Wielandt inequality. Part~(b) is the reason for working with $W_1$: a change of state by a unitary costs
the state commutator linearly.

\subsubsection*{The profile of \texorpdfstring{$K$}{K}}

\begin{theorem}[The profile of $c$ in a state]\label{thm:K-profile}
Let $\phi$ assign unitaries $A,B,C,X,Y$ on $\C^d$ to $a,b,c,x,y$, and let $\rho$ be a density matrix on $\C^d$ with $\xi=\rho^{1/2}$.
Suppose that, for some $e>0$,
\begin{enumerate}[label=(\roman*),itemsep=1pt]
\item every relator $r$ in~\eqref{eq:K} has $\hs{(\phi(r)-I)\xi}\le e$;
\item $\hs{[\phi(s),\xi]}\le e$ for $s\in\{a,b,c,x,y\}$;
\item $\hs{(C+I)\xi}\le e$.
\end{enumerate}
Then $d\ge1/(29e)$. In the notation of \cref{def:state-profile}, $\hlps_K(c;e)\ge1/(29e)$.
\end{theorem}

The two relations are read off as near-invariances of $\mu=\mu^\rho_Y$, each up to a coupling in the state (\cref{lem:coupling}(a)) and a
change of state by a unitary (\cref{lem:coupling}(b)); \cref{lem:circle} then compares $\mu$ with Haar measure. That comparison does not
involve the dimension, and we state it on its own.

\begin{proposition}[The spectral measure of $y$ in a state]\label{prop:K-W1}
Under the hypotheses of \cref{thm:K-profile}, the spectral measure of $Y=\phi(y)$ in $\rho$ satisfies
\[
W_1(\mu^\rho_Y,\mathsf h)\ \le\ \frac{29\pi}2\,e .
\]
\end{proposition}

The bound uses all three hypotheses: the state commutators~(ii) pay for the changes of state in \cref{lem:coupling}(b), and~(iii), which
says that $c$ acts as $-1$ in the state, turns the relations $yay^{-1}=b$ and $c=ab$ into the half-turn.

\begin{proof}
Let $\mu=\mu^\rho_Y$, $\delta_{\mathsf D}=W_1(\mathsf D_*\mu,\mu)$ and $\delta_{\mathsf S}=W_1(\mathsf S_*\mu,\mu)$.

\emph{Doubling.} Put $P=Y^2$ and $Q=XYX^*$. The relator $r=xyx^{-1}y^{-2}$ has $\phi(r)=QP^*$, so $(Q-P)\xi=(\phi(r)-I)P\xi=(\phi(r)-I)\xi
P+(\phi(r)-I)[P,\xi]$, and $\hs{[P,\xi]}\le2e$; hence $\hs{(Q-P)\xi}\le5e$. Also $\hs{[Q,\xi]}\le3e$, so $\hs{P\xi-\xi Q}\le\hs{(P-Q)\xi}+\hs{[Q,\xi]}\le8e$.
By \cref{lem:coupling}(a) and~\eqref{eq:rates-spectral}, $W_1(\mathsf D_*\mu,\mu^{X^*\rho X}_Y)\le4\pi e$, and by \cref{lem:coupling}(b),
$W_1(\mu^{X^*\rho X}_Y,\mu)\le\pi e$. So $\delta_{\mathsf D}\le5\pi e$.

\emph{Half-turn.} The word $w=a^{-1}yay^{-1}c^{-1}$ is freely equal to
\[
\bigl(a^{-1}(yay^{-1}b^{-1})a\bigr)\,(a^2)^{-1}\,\bigl(c(c^{-1}ab)c^{-1}\bigr),
\]
a product of three conjugates of relators with conjugators of length $1$, $0$ and $1$, so $\hs{(\phi(w)-I)\xi}\le3e+e+3e=7e$ by
\cref{lem:calculus}(d). Since $A^*YA=\phi(w)CY$,
\[
\hs{(A^*YA-CY)\xi}=\hs{(\phi(w)-I)CY\xi}\le7e+2\hs{[CY,\xi]}\le11e,
\]
and $\hs{(CY+Y)\xi}=\hs{(C+I)Y\xi}\le\hs{(C+I)\xi}+2\hs{[Y,\xi]}\le3e$. So $\hs{(A^*YA+Y)\xi}\le14e$ and
$\hs{(-Y)\xi-\xi(A^*YA)}\le14e+\hs{[A^*YA,\xi]}\le17e$. By \cref{lem:coupling}(a) and~\eqref{eq:rates-spectral},
$W_1(\mathsf S_*\mu,\mu^{A\rho A^*}_Y)\le\frac{17}2\pi e$, and by \cref{lem:coupling}(b) with $U=A^*$, $W_1(\mu^{A\rho A^*}_Y,\mu)\le\pi e$.
So $\delta_{\mathsf S}\le\frac{19}2\pi e$.

\emph{Conclusion.} By \cref{lem:circle}(c), $W_1(\mu,\mathsf h)\le\delta_{\mathsf D}+\delta_{\mathsf S}\le\bigl(5+\frac{19}2\bigr)\pi
e=\frac{29}2\pi e$.
\end{proof}

\begin{proof}[Proof of \cref{thm:K-profile}]
The measure $\mu^\rho_Y$ is supported on the angles of the at most $d$ eigenvalues of $Y$, so by \cref{lem:circle}(d) and
\cref{prop:K-W1}, $\frac\pi{2d}\le W_1(\mu^\rho_Y,\mathsf h)\le\frac{29}2\pi e$, that is $d\ge1/(29e)$. Conditions (i)--(iii) are
\ref{prelim:S1}--\ref{prelim:S3} for $K$ and $w=c$, which gives the second sentence.
\end{proof}

For $\rho=I/d$ the state commutators vanish, and (i)--(iii) say that $\phi$ is an $e$-representation with $\fn{\phi(c)+I}\le e$. Together
with an explicit model this pins down the profile of $c$.

\begin{example}[A model]\label{ex:K-model}
Let $d=2m$, $z_j=e^{2\pi ij/d}$ for $0\le j<d$, and on $\C^d$ put $Y=\operatorname{diag}(z_j)$, let $A$ exchange $e_j$ and $e_{j+m}$
(indices mod $d$), $B=-A$, $C=-I$, and let $X$ be the permutation matrix with $Xe_{\varpi(j)}=e_j$, where $\varpi(j)=2(j\bmod m)+\lfloor
j/m\rfloor$. Then $A^2=B^2=C^2=I$ and $C^{-1}AB=A^2=I$; $YAY^*=-A=B$ and $YBY^*=A$, because $z_{j+m}/z_j=-1$; and $C$ is central. Finally
$XYX^*=\operatorname{diag}(z_{\varpi(j)})$, so $XYX^*Y^{-2}$ is $I$ on the first $m$ coordinates and $e^{2\pi i/d}I$ on the last $m$. With
the flat state, every relator of~\eqref{eq:K} but $xyx^{-1}y^{-2}$ has defect $0$, that one has defect
$\fn{XYX^*Y^{-2}-I}=\sqrt2\sin(\pi/d)$, and (ii), (iii) hold with $0$. This is a variant of the model of~\cite[Lemma~3.8]{SV18}.
\end{example}

\begin{corollary}[The profile of $c$]\label{cor:K-profile}
For every $e>0$,
\[
\frac1{29e}\ \le\ \hlps_K(c;e)\ \le\ \hlp_K(c;2,e)\ \le\ 2\Bigl\lceil\frac\pi{\sqrt2\,e}\Bigr\rceil .
\]
For each $0<\delta\le2$ there are constants $c_\delta,C_\delta>0$ with $c_\delta/\eps\le\hlp_K(c;\delta,\eps)\le C_\delta/\eps$ for
$0<\eps\le2$.
\end{corollary}

\begin{proof}
The first inequality is \cref{thm:K-profile}. The second holds for every presentation: an $e$-representation with $\fn{\phi(c)-I}\ge2$ has
$\phi(c)=-I$ by~\eqref{eq:unitary-distance}, and with the flat state it satisfies (S1)--(S3). For the third, \cref{ex:K-model} with
$m=\lceil\pi/(\sqrt2e)\rceil$ has $\phi(c)=-I$ and defects at most $\sqrt2\sin(\pi/(2m))\le\sqrt2\pi/(2m)\le e$. For $\delta<2$,
$\hlp_K(c;2,C\eps)\le\hlp_K(c;\delta,\eps)\le\hlp_K(c;2,\eps)$ for a constant $C\ge1$ depending on $\delta$~\cite[Proposition~2.7]{SV18},
since $c$ is a central involution represented by a generator. So $\hlp_K(c;\delta,\eps)\ge1/(29C\eps)$, and for $\eps\le2$,
$\hlp_K(c;\delta,\eps)\le2(\pi/(\sqrt2\eps)+1)\le(\sqrt2\pi+4)/\eps$.
\end{proof}

\Cref{cor:K-profile} improves the lower bound $C\eps^{-2/3}$ of~\cite[Lemma~3.5 and Proposition~3.2]{SV18} to the order of their upper
bound. Their argument passes from the normalized norm to the operator norm, which costs a factor $\sqrt d$; the argument above never uses
the operator norm.

\subsubsection*{The game}

Let $\Gsv=G(A,b)$ be the linear system game of~\cite[Theorem~1.1]{SV18}, with $m$ equations, $n$ variables, $s=\max_i\lvert V_i\rvert$ and
solution group $\Gamma=\Gamma(A,b)$, the solution group of~\cite[Proposition~6.1]{SV18}. Its question distribution is that
of~\cite[Section~4]{SV18}: the equation $i$ and the variable $j$ are independent and uniform, and a pair with $j\notin V_i$ is always won;
so $\pi_{\min}=1/(mn)$ (\cref{def:lcs}). The construction~\cite[Propositions~6.1 and~6.2]{SV18} applies Slofstra's embedding
theorem~\cite{Slo20}, in the form of~\cite[\S4]{Slo19}, to $K$, with distinguished central involution $c$. It supplies a homomorphism
$\Psi:F(a,b,c,x,y)\to F(x_1,\dots,x_n,J)$ that descends to an embedding of $K$ into $\Gamma$, and a word $z$ in $x_1,\dots,x_n,J$, the
image of an added involution $t$ with $tct=cJ$, such that $[\Psi(c),z]=J$ in $\Gamma$~\cite[proof of Proposition~6.2]{SV18}. (The proof
there prints the relation as $twt=J$, for the distinguished involution $w$; the relation used is $twt=wJ$, that of~\cite[Lemma~5.2]{Slo19},
which is what gives $[w,t]=J$.) For the lower bound we use only that $\Psi$ descends to a homomorphism $K\to\Gamma$ and this identity.
Since $J$ is nontrivial in approximate representations of $\Gamma$~\cite[Proposition~6.1]{SV18}, $\omq(\Gsv)=1$ (see the discussion after
\cref{def:lcs}), so $\Ent(\Gsv,\eps)$ counts strategies that win with probability at least $1-\eps$. Qualitatively, this nontriviality comes
from Slofstra's embeddings and his tensor-power argument~\cite[Lemma~5.2]{Slo19}; \cref{lem:sv-lift} below gives it again, with rates.

The upper bound $\Ent(\Gsv,\eps)\le C'\eps^{-1/2}$ of~\cite[Theorem~1.1]{SV18} rests on the second half of the proof
of~\cite[Proposition~6.2]{SV18}. That argument takes from~\cite[Proposition~4.8]{Slo19} a $2d$-dimensional lift of a $d$-dimensional
approximate representation of $K$ with $c\mapsto(-I_d)\oplus I_d$, and then lets $t$ exchange the two blocks. But~\cite[Proposition~4.8]{Slo19}
removes the non-involutive generators one at a time, and each removal doubles the dimension and sends every involutive generator $s$ to
the image of $s$ plus an identity block. The group $K$ has two non-involutive generators, $y<x$, and no presentation of $K$ in
the sense of~\cite[\S6]{SV18} has fewer: the quotient of $K$ by $a,b,c$ is $\langle x,y\mid xyx^{-1}=y^2\rangle$, which is torsion-free and
not cyclic, so it is also the quotient of $K$ by all its involutions, and that quotient is generated by the non-involutive generators of
any such presentation. So the two lifts give $c\mapsto W=(-I_d)\oplus I_{3d}$. For every unitary $T$ on $\C^{4d}$ the $(+1)$-eigenspaces of $W$ and $TWT^*$ have dimension $3d$, so they
meet in dimension at least $2d$, where $TWT^*+W=2I$; hence $\fn{TWT^*+W}^2\ge2$, and the relation $tct=cJ$ with $J\mapsto-I$ fails by at
least $\sqrt2$, whatever the precision of the representation. So for $K$ this last step, as written in~\cite{SV18} (arXiv:1711.10676v2), does not produce the
representation it needs; with at most one non-involutive generator it is complete, and that is the case of \cref{thm:embedding}. This
affects only the upper bound $\hlp(J;2,\eps)\le N\,\hlp(c;2,C'\eps)$ of~\cite[Prop.~6.2]{SV18}. Its lower bound, the nontriviality of $J$
and the value one of the game do not use this step: they use only that the embedding of~\cite[Prop.~4.8]{Slo19} keeps $c$ nontrivial in
approximate representations, then~\cite[Lemma~5.2 and Prop.~4.2]{Slo19} and, for the value, \cite[Cor.~5.3]{SV18}. For $K$ the upper
bound itself holds: \cref{lem:sv-lift} below supplies a balanced lift for the explicit models of $K$, and with it the upper bounds
of~\cite[Prop.~6.1 and Thm.~1.1]{SV18} hold as stated. For the models of \cref{ex:K-model}, in which $X$ is a permutation matrix, the second removal can be done without
doubling, and the marked involution stays balanced.

\begin{lemma}[A balanced lift for $K$]\label{lem:sv-lift}
Let $m\ge1$, $d=2m$, and let $\phi$ be the model of \cref{ex:K-model} on $\C^d$, whose relators have defect at most
$e_m=\sqrt2\sin(\pi/(2m))$. Then the solution group $\Gamma$ of $\Gsv$ has a $C_\Gamma e_m$-representation $\gamma$ on $\C^{8d}$ with
$\gamma(J)=-I$, where $C_\Gamma\le75$ is the constant of~\cite[Remark~4.5]{Slo19}.
\end{lemma}

\begin{proof}
We follow the construction of $\Gamma$ in~\cite[proof of Proposition~6.2]{SV18} and lift $\phi$ through each of its steps. Write
$A,B,C,X,Y$ for the images of $a,b,c,x,y$, and use the notation $\Omega$, $U\oplus V$ and $\operatorname{adiag}(U,V)$ of
\cref{gadget:conv}. As an extended presentation in the sense of~\cite[\S6]{SV18}, $K$ has the involutive generators $a,b,c$, the linear
set $\{a,b,c\}$ with its commutations, which Slofstra and Vidick add and the model satisfies exactly, the conjugacy relations $yay^{-1}=b$,
$yby^{-1}=a$ and $xcx^{-1}=c$, and the non-involutive generators $y<x$ with the relation $xyx^{-1}=y^2$. In the model only the relator
$r=xyx^{-1}y^{-2}$ has nonzero defect, and that defect is $e_m$.

\emph{Step 1: removing $y$}~\cite[Proposition~4.8]{Slo19}. Add involutions $f,v$, with $y=fv$ in mind, and the relation that $f$ commutes
with $x$. Replace $yay^{-1}=b$ by $vav=Z_a$ and $fZ_af=b$, and $yby^{-1}=a$ by $vbv=Z_b$ and $fZ_bf=a$, with ancilla involutions
$Z_a,Z_b$. Once $y=fv$ and $fx=xf$, the relation $xyx^{-1}=y^2$ reads $xvx^{-1}=vfv$; replace it by $vfv=W_0$ and $xvx^{-1}=W_0$, with an
ancilla involution $W_0$. Then $y$ occurs only in $y=fv$, and is deleted. On $\C^d\oplus\C^d$ put
\[
\gamma_1(s)=\phi(s)\oplus I\ \ (s=a,b,c),\qquad\gamma_1(f)=\Omega,\qquad\gamma_1(v)=\operatorname{adiag}(Y^*,Y),\qquad\gamma_1(x)=X\oplus X,
\]
and give $Z_a,Z_b,W_0$ their defining values $\gamma_1(vav)$, $\gamma_1(vbv)$, $\gamma_1(vfv)$. This is the lift
of~\cite[Proposition~4.8]{Slo19}, with $\gamma_1(f)\gamma_1(v)=Y\oplus Y^*$. The relations that define the ancillas hold exactly,
and so do $f^2=v^2=e$ and $fx=xf$. Every other relator is mapped to a block-diagonal matrix each of whose two blocks is the identity or a conjugate
of $\phi(r')^{\pm1}$ for a relator $r'$ of $K$: for instance
$\gamma_1(fZ_afb^{-1})=\phi(yay^{-1}b^{-1})\oplus I$, and $\gamma_1(xvx^{-1}W_0^{-1})$ is conjugate to $\phi(r)\oplus Y^{-2}\phi(r)^{-1}Y^2$. So
$\gamma_1$ is an $e_m$-representation of the presentation reached, and $\gamma_1(c)=(-I_d)\oplus I_d$.

\emph{Step 2: removing $x$ without doubling.} Now $x$ is the only non-involutive generator. The construction
of~\cite[Proposition~4.8]{Slo19} adds involutions
$f',v'$, with $x=f'v'$ in mind, and replaces each of the three relations $xsx^{-1}=s'$, namely $(s,s')=(c,c)$, $(f,f)$ and $(v,W_0)$, by
$v'sv'=Z'_s$ and $f'Z'_sf'=s'$, with an ancilla involution $Z'_s$; it adds no commutation relations, since no non-involutive generator
follows $x$, and then deletes $x$. The lift used there would double the dimension again. Instead we use that $X$ is a permutation matrix.
Reflecting each cycle of the permutation, $a_i\mapsto a_{-i}$ on a cycle $(a_0\,a_1\,\cdots\,a_{k-1})$, gives an involutive permutation
matrix $\tau$ with $\tau X\tau=X^{-1}$. Put
\[
\gamma_2(f')=\tau\oplus\tau,\qquad\gamma_2(v')=\tau X\oplus\tau X,\qquad\gamma_2(Z'_s)=\gamma_2(v')\gamma_1(s)\gamma_2(v'),
\]
and $\gamma_2(s)=\gamma_1(s)$ for the other generators. The matrix $\tau X$ is real, and $(\tau X)^*=X^{-1}\tau=\tau X$ and $(\tau X)^2=X^{-1}X=I$,
so $\gamma_2(f')$ and $\gamma_2(v')$ are involutions, and $\gamma_2(f')\gamma_2(v')=X\oplus X=\gamma_1(x)$. Hence
$\gamma_2(f'Z'_sf's'^{-1})=\gamma_1(xsx^{-1}s'^{-1})$: every relation keeps its residual. The relations $v'sv'=Z'_s$ hold exactly, and
${Z'_s}^2$ is conjugate to $s^2$. So $\gamma_2$ is an $e_m$-representation, on $\C^{2d}$, of a presentation all of whose generators are
involutions, and still $\gamma_2(c)=(-I_d)\oplus I_d$.

\emph{Step 3: the commutator}~\cite[proof of Proposition~6.2]{SV18}. Adjoin the central involution $J$ and an involution $t$ with
$tct=cJ$. In linear-plus-conjugacy form this takes one more involution $u$, with the linear relation $cu=J$ and the conjugacy relation
$tct=u$, as in Step~2 of the proof of \cref{lem:gadget}; Slofstra and Vidick leave this encoding implicit, and any encoding through
auxiliary involutions that are determined by $t$, $c$ and $J$ works the same way. Put $\gamma_3(J)=-I$, let $\gamma_3(t)=\Omega$ exchange the
two copies of $\C^d$, and put $\gamma_3(u)=-\gamma_2(c)=I_d\oplus(-I_d)$. Because $\gamma_2(c)=(-I_d)\oplus I_d$ has its two eigenvalues with
equal multiplicities, $\gamma_3(t)\gamma_3(c)\gamma_3(t)=\gamma_3(u)$; also $\gamma_3(c)\gamma_3(u)=-I=\gamma_3(J)$, and $-I$ commutes with
everything. So all the new relations hold exactly, and $\gamma_3$ is an $e_m$-representation of the resulting linear-plus-conjugacy
presentation on $\C^{2d}$ with $J\mapsto-I$.

\emph{Step 4: the solution group}~\cite[Lemma~4.4 and Proposition~4.2]{Slo19}. These two steps turn the linear-plus-conjugacy presentation
into the solution group $\Gamma$. By~\cite[Remark~4.5]{Slo19} they turn every $2d$-dimensional $e$-representation of it into an
$8d$-dimensional $C_\Gamma e$-representation of $\Gamma$ that restricts to four copies of the given one on its generators, with
$C_\Gamma\le75$. Applied to $\gamma_3$, this gives $\gamma$, and $\gamma(J)=-I$.
\end{proof}

The program \texttt{anc/sv\_balanced\_lift.py} builds these four steps for $m\le8$ and evaluates every relator. Steps~1--3 keep every defect at
most $e_m$, and all $2373$ relators of the solution group so obtained, with the products of the equations taken in every order, have
defect at most $e_m/\sqrt2$, far below $75e_m$. As a negative control it builds the two literal lifts and confirms the bound
$\fn{TWT^*+W}^2\ge2$, which a block permutation attains (\cref{cert:what}).

\begin{theorem}[Entanglement of the Slofstra--Vidick game]\label{thm:sv-game}
There are constants $c_0,C_0>0$ such that $c_0\,\eps^{-1/2}\le\Ent(\Gsv,\eps)\le C_0\,\eps^{-1/2}$ for $0<\eps\le1$.
\end{theorem}

For the upper bound, \cref{lem:sv-lift} and the return direction with the flat state give strategies on a maximally entangled state of
local dimension $16m$ that lose with probability $O(1/m^2)$. For the lower bound, a near-optimal strategy gives, by \cref{thm:transfer}, a
model of $\Gamma$ in Bob's reduced state at precision $O(\sqrt\eps)$, in which $J=-I$. Composed with $\Psi$ it is a model of $K$ in which
$c$ is visible, by the identity $[\Psi(c),z]=J$; on the sector where the rounding of $\phi(\Psi(c))$ is $-1$ it satisfies the hypotheses
of \cref{thm:K-profile}, which bounds the Schmidt rank by $\Omega(\eps^{-1/2})$.

\begin{proof}
\emph{The upper bound.} Let $C_s$ be the constant of \cref{thm:return} and $C_\Gamma$ that of \cref{lem:sv-lift}. For $m\ge1$,
\cref{lem:sv-lift} gives a $C_\Gamma e_m$-representation of $\Gamma$ on $\C^{16m}$ with $J\mapsto-I$, where
$e_m=\sqrt2\sin(\pi/(2m))\le\pi/(\sqrt2\,m)$. By \cref{thm:return} with the flat state, the strategy built there on $\Phi_{16m}$ loses every
question pair, hence a random one, with probability at most $C_sC_\Gamma^2e_m^2\le\pi^2C_sC_\Gamma^2/(2m^2)$. Given $0<\eps\le1$, take
$m=\lceil\pi C_\Gamma\sqrt{C_s/(2\eps)}\,\rceil$. Then the losing probability is at most $\eps$, and since $\eps\le1$,
\[
\Ent(\Gsv,\eps)\ \le\ \Entme(\Gsv,\eps)\ \le\ 16m\ \le\ 16\Bigl(\pi C_\Gamma\sqrt{C_s/2}+1\Bigr)\eps^{-1/2}.
\]
This is the upper bound of~\cite[Theorem~1.1]{SV18}, with $C_0=16(\pi C_\Gamma\sqrt{C_s/2}+1)$; it also shows $\omq(\Gsv)=1$ once more.

\emph{The lower bound.} Take a strategy that wins with probability at least $1-\eps$, with
Schmidt rank $r$, and let $\phi$ and $\rho$ on $\C^r$ be as in \cref{thm:transfer}, with $\theta=2\sqrt{mn\eps}$ and $e_0=\max(s,4)\theta$.
Put $\phi_K=\phi\circ\Psi$ and $W=\phi_K(c)$.

\emph{Constants.} The following words are trivial in $\Gamma$: $\Psi(r)$ for each relator $r$ of~\eqref{eq:K}; $\Psi([u,c])$ for
$u\in\{a,b,c,x,y\}$, as $c$ is central in $K$; and $[\Psi(c),z]J^{-1}$. Fix for each of them an expression as a product of conjugates of
relators of $\Gamma$. These finitely many expressions do not depend on $\eps$ or on the strategy. By \cref{lem:calculus}(d),(a) and
\cref{thm:transfer}, which bounds the defects of the relators of $\Gamma$ by $e_0$ and the state commutators of the letters $x_j$ by
$\sqrt2\theta$ (and $\phi(J)=-I$ commutes with $\xi$), there is a constant $C_1$, depending only on $\Gsv$ and the chosen expressions, such that each of these words has
defect $\hs{(\phi(\cdot)-I)\xi}\le C_1\theta$, and $\hs{[\phi_K(u),\xi]},\ \hs{[\phi(z),\xi]}\le C_1\theta$ for $u\in\{a,b,c,x,y\}$.

\emph{Visibility.} Let $V=\phi(z)$ and $\tau=\hs{(W-I)\xi}$. Since $\phi(J)=-I$, $\phi([\Psi(c),z]J^{-1})=-WVW^*V^*$, so
$\hs{(WVW^*V^*+I)\xi}\le C_1\theta$. From $WVW^*V^*-I=(W-I)+WV(W^*-I)V^*$,
\[
\hs{(WVW^*V^*-I)\xi}\le\tau+\hs{(W^*-I)(\xi V^*+[V^*,\xi])}\le2\tau+2\hs{[V,\xi]}\le2\tau+2C_1\theta .
\]
As $2=\hs{2\xi}$, we get $\tau\ge1-\frac32C_1\theta$.

\emph{The sector.} Let $J'=\mathrm r(W)$. By \cref{lem:rounding} and the relator $c^2$, $\hs{(J'-W)\xi}\le C_1\theta$ and
$\hs{[J',\xi]}\le3C_1\theta$. Let $P_-=\frac12(I-J')$ and $p_-=\Tr(P_-\rho)=\frac14\hs{(I-J')\xi}^2$. Then $\hs{(I-J')\xi}\ge\tau-C_1\theta\ge1-\frac52C_1\theta$,
so $p_-\ge\frac1{16}$ when $\theta\le1/(5C_1)$; assume this. For $u\in\{a,b,c,x,y\}$ and $U=\phi_K(u)$, the group commutator
$UWU^*W^*=\phi(\Psi([u,c]))$ has defect at most $C_1\theta$, so $\hs{[U,W]\xi}=\hs{(UWU^*W^*-I)WU\xi}\le C_1\theta+4C_1\theta$. With
$N=J'-W$, $\op N\le2$ and $\hs{[U,N]\xi}\le\hs{N\xi}+\hs{N\xi}+\op N\hs{[U,\xi]}\le4C_1\theta$, so $\hs{[U,J']\xi}\le9C_1\theta$.

\emph{The profile in the sector.} Let $\rho_-$ and $\xi_-$ be as in \cref{lem:sector-state}, with $\sqrt{p_-}\ge\frac14$. By
\cref{lem:sector-state}(c), $\hs{[U,\xi_-]}\le4(9+3+3)C_1\theta=60C_1\theta$ for $U=\phi_K(u)$. By \cref{lem:sector-state}(b) with
$\op{\phi_K(r)-I}\le2$, every relator of~\eqref{eq:K} has $\hs{(\phi_K(r)-I)\xi_-}\le4(C_1+6C_1)\theta=28C_1\theta$. Finally $(J'+I)P_-=0$,
so $(J'+I)\xi_-=(J'+I)(\xi_--P_-\xi P_-/\sqrt{p_-})$, of norm at most $\hs{[J',\xi]}/\sqrt{p_-}$ by \cref{lem:sector-state}(a), and
\[
\hs{(W+I)\xi_-}\le\hs{(W-J')\xi_-}+\hs{(J'+I)\xi_-}\le4(C_1+6C_1)\theta+12C_1\theta=40C_1\theta .
\]
So $\phi_K$ and $\rho_-$ satisfy (i)--(iii) of \cref{thm:K-profile} on $\C^r$ with $e=60C_1\theta$, and
$r\ge1/(1740\,C_1\theta)=\eps^{-1/2}/(3480\,C_1\sqrt{mn})$.

If $\theta>1/(5C_1)$, that is $\eps>\eps_0=1/(100C_1^2mn)$, then $\Ent(\Gsv,\eps)\ge1\ge\sqrt{\eps_0}\,\eps^{-1/2}$. So the lower bound
holds for all $\eps>0$ with $c_0=\min\bigl(1/(3480C_1\sqrt{mn}),\sqrt{\eps_0}\bigr)$, since a strategy of local dimension $d$ has Schmidt rank
at most $d$.
\end{proof}

\begin{remark}\label{rates:rem-sv}
The lower bound $C\eps^{-1/6}$ of~\cite[Theorem~1.1]{SV18} passes from the state to a spectral projection of it, and from the
normalized norm to the operator norm; the proof above keeps the state throughout. The bound is on the Schmidt rank, and so on the local
dimension: $\log\Ent(\Gsv,\eps)=\frac12\log(1/\eps)+O(1)$. It is not a bound on the entanglement entropy, since \cref{thm:K-profile} counts
the atoms of a spectral measure, whose masses may be unequal.
\end{remark}

\emph{Theorem~F is \cref{cor:K-profile} with \cref{thm:sv-game}.}
